% !PRINTED 385
\weberpage{385}
%<<CONT
aufgeht. In diesem Falle ist aber $(D, p^s) = (\Delta, p^s)$ [\S\,85 (9)],
und es ist also zu beweisen:
\[
\chi(\delta, p^k) \sum_{0,k}^{s} (\Delta, p^s) = \sum_{0,k}^{s} (\delta, p^s)(\delta', p^{k-s}).
\]
Geht nun $p$ nicht in $\delta$ auf, so ist
\[
\chi(\delta, p^k) = (\delta, p^k),
\]
und wenn man rechts jedes Glied unter dem Summenzeichen mit
$(\delta, p^{k-s})^2 = 1$ multipliziert, so erhält man übereinstimmend mit
der linken Seite
\[
(\delta, p^k) \overset{s}{\sum} (\Delta, p^{k-s}). % [sic: lower limit omitted in source]
\]
Geht aber $p$ in $\delta$ und folglich nicht in $\delta'$ auf, so ist $\chi(\delta, p^k)$
$= (\delta', p^k)$, multipliziert man also ebenso mit $(\delta', p^s)^2$, so folgt
das gleiche, wodurch (22) bewiesen ist.

8. \gesperrt{Ist $\chi(\delta, n)$ von $\chi(\delta', n)$ verschieden, so ist}
\weq{23}{\Sigma(D, \bar e) = 0, \quad \Sigma(\delta, e)(\delta', e') = 0.}
Denn da sich die beiden Summen nicht ändern, wenn $\delta$ mit
$\delta'$ vertauscht wird, so folgt dies aus (22).

9. \gesperrt{Sind $\delta_1, \delta_2$ zwei Stammteiler von $D$ und $\delta_1\delta_2$ ihr
symbolisches Produkt, so ist}
\weq{24}{\chi(\delta_1, n)\, \chi(\delta_2, n) = \chi(\delta_1\delta_2, n).}
Diese Formel braucht wegen (21) nur für den Fall bewiesen
zu werden, daß $n$ eine Primzahl ist.

Es ist aber symbolisch
\[
(\delta_1\delta_2)' = \delta_1'\delta_2 = \delta_1\delta_2',
\]
und hiernach ergibt sich (24) leicht aus der Definition (19).

\begin{center}
§ 105.\quad \gesperrt{Anwendung des Legendreschen Symbols.}
\end{center}

Will man sich der bekannteren Bezeichnung nach \gesperrt{Legendre}
und \gesperrt{Jacobi} bedienen, so gestaltet sich die Sache folgendermaßen.
Bedeutet $l$ eine ungerade Primzahl und $a$ eine durch $l$ nicht
teilbare Zahl, so ist nach \gesperrt{Legendre}:
\weq{1}{\left(\frac{a}{l}\right) = +1,}
\gesperrt{wenn $a$ quadratischer Rest von $l$ ist,}
\[
\left(\frac{a}{l}\right) = -1,
\]
%CONT>>
